% =====================================================================
%  Determinantes.tex
%  ========================
%
%
%  Descripción:
%  ------------
%  Presentación en beamer sobre determinantes y los tres métodos para
%  calcularlos: eliminación gaussiana, la fórmula grande sobre
%  permutaciones y la expansión por cofactores. Sigue el capítulo 5
%  de Strang, "Introduction to Linear Algebra", secciones 5.1 y 5.2,
%  que es el extracto que se usó como fuente. Las láminas llevan
%  solo lo esencial; lo que se dice en voz alta vive en
%  .GuionDeterminantes.txt, lámina por lámina.
%  Metadata:
%  ----------
%   * Author : zxxz6 (Bryan Violante Arriaga)
%   * Version: 4.6.0
%
%
%  History:
%  ------------
%  Author      Date            Description
%  zxxz6       12/09/2026      Creation
%
%
% =====================================================================
\documentclass[aspectratio=169,11pt]{beamer}

% =====================================================================
%  DATOS DE LA PRESENTACIÓN
% =====================================================================
\newcommand{\materia}{Matemáticas para Computación}
\newcommand{\periodo}{Otoño 2026}
\newcommand{\alumnouno}{Bryan Ezequiel Violante Arriaga}
\newcommand{\alumnodos}{Violeta Del Pilar Lagunes Viveros}

% ==================== IDIOMA Y CODIFICACIÓN ====================
\usepackage[utf8]{inputenc}
\usepackage[T1]{fontenc}
\usepackage[spanish,es-tabla,es-noshorthands]{babel}
\usepackage{lmodern}

% ==================== MATEMÁTICAS Y TABLAS ====================
\usepackage{amsmath,amssymb,mathtools}
\usepackage{booktabs}
\usepackage{array}
\usepackage{xcolor}
\usepackage{tikz}
\usetikzlibrary{matrix,calc,backgrounds,babel}

% ==================== TEMA ====================
\usetheme{metropolis}
\metroset{block=fill,numbering=none,progressbar=frametitle}

% --- la misma paleta de las libretas ---
\definecolor{acento}{HTML}{1F4E79}
\definecolor{fondoIdea}{HTML}{EAF2FB}
\definecolor{fondoOjo}{HTML}{FDF0E6}
\setbeamercolor{palette primary}{bg=acento,fg=white}
\setbeamercolor{frametitle}{bg=acento,fg=white}
\setbeamercolor{progress bar}{fg=orange!70!black}
\setbeamercolor{alerted text}{fg=orange!75!black}

% --- las dos cajas que se usan a lo largo del deck ---
\setbeamercolor{block title}{bg=acento!15,fg=acento}
\setbeamercolor{block body}{bg=acento!7}
\setbeamercolor{block title alerted}{bg=fondoOjo,fg=orange!45!black}
\setbeamercolor{block body alerted}{bg=fondoOjo!60}

% --- abreviaturas que se repiten en cada lámina ---
\newcommand{\dt}[1]{\det #1}
\newcommand{\abs}[1]{\left\lvert #1 \right\rvert}
% El determinante escrito con barras, que es como aparece en Strang
\newcommand{\dbar}[1]{%
  \begin{vmatrix} #1 \end{vmatrix}%
}
\newcommand{\mat}[1]{%
  \begin{bmatrix} #1 \end{bmatrix}%
}

% --- entrada tachada: se atenua, no se borra, para que se vea
%     de donde sale el menor ---
\newcommand{\ta}[1]{{\color{gray!40}#1}}

% --- remate de lamina: una frase, sin barra de titulo vacia ---
\newcommand{\remate}[1]{%
  \vspace{0.3em}
  \begin{center}\itshape\small #1\end{center}%
}

\title{Determinantes}
\author{\alumnouno \and \alumnodos}
\institute{\materia, INAOE}
\date{\periodo}

% =====================================================================
\begin{document}
% =====================================================================

\maketitle

\begin{frame}{Disclaimer :P}

    El material se preparó usando inteligencia artificial como
    herramienta de apoyo, basado en fuentes verificables y confiables (chatgepete)

\end{frame}

\begin{frame}[t]{Contenido}
  \small
  \tableofcontents
\end{frame}

% =====================================================================
%  1. Qué es un determinante
% =====================================================================
\section{Qué es un determinante}

\begin{frame}{Definicion}
  \begin{block}{Matriz cuadrada $\longrightarrow$ un número}
    \[
      A^{-1} \text{ existe}
      \iff
      \dt{A} \neq 0
    \]
  \end{block}

  \begin{itemize}
    \item Notación: $\dt{A}$ o $\abs{A}$.
    \item Las barras son de la matriz, no valor absoluto:
          $\abs{A}$ puede ser negativo.
    \item $\dt{(A^{-1})} = 1/\dt{A}$.
  \end{itemize}
\end{frame}

\begin{frame}{Det de Matriz $2 \times 2$}
  \[
    A = \mat{a & b \\ c & d}
    \qquad
    \dt{A} = ad - bc
    \qquad
  \]

  \begin{itemize}
    \item Para $n$ grande no sirve. Explica, no calcula.
  \end{itemize}
\end{frame}


\begin{frame}{Tres formas de calcularlo}
  \begin{enumerate}
    \item \textbf{Eliminación gaussiana.} Triangular, multiplicar
          pivotes. $\sim n^3/3$.
    \item \textbf{Fórmula grande.} Suma de $n!$ términos.
    \item \textbf{Cofactores.} Orden $n$ a orden $n-1$.
  \end{enumerate}

  \remate{Mismo número, distinto costo.}
\end{frame}

% =====================================================================
%  2. Las diez propiedades
% =====================================================================
\section{Las propiedades}

\begin{frame}{Regla 1: Determinante de la matriz identidad}
  \[
    \dt{I} = 1
  \]

  \begin{block}{Ejemplo}
    \[
      \dbar{1 & 0 \\ 0 & 1} = 1
      \qquad
      \dbar{1 & 0 & 0 \\ 0 & 1 & 0 \\ 0 & 0 & 1} = 1
    \]
  \end{block}

  \remate{Es la normalización. Sin ella todo se podría
          escalar por una constante.}
\end{frame}

\begin{frame}{Regla 2: intercambiar renglones cambia el signo}
  \begin{block}{Ejemplo}
    \[
      \dbar{1 & 2 \\ 3 & 4} = 4 - 6 = -2
      \qquad
      \dbar{3 & 4 \\ 1 & 2} = 6 - 4 = +2
    \]
  \end{block}

  \begin{itemize}
    \item $P$ es $I$ con los renglones reordenados.
    \item $\dt{P} = +1$ con un número par de intercambios,
          $\dt{P} = -1$ con uno impar.
  \end{itemize}
\end{frame}

\begin{frame}{Regla 3: lineal en cada renglón}
  \begin{block}{Escalar un renglón}
    \[
      \dbar{ta & tb \\ c & d} = t\,\dbar{a & b \\ c & d}
    \]
  \end{block}

  \begin{block}{Sumar en un renglón}
    \[
      \dbar{a + a' & b + b' \\ c & d}
      = \dbar{a & b \\ c & d} + \dbar{a' & b' \\ c & d}
    \]
  \end{block}

  \begin{alertblock}{Un renglón a la vez}
    El segundo renglón no se movió en ninguna de las dos.
  \end{alertblock}
\end{frame}

\begin{frame}{Regla 4: dos renglones iguales dan cero}
  \begin{block}{Ejemplo}
    \[
      \dbar{2 & 7 \\ 2 & 7} = 14 - 14 = 0
      \qquad
      \dbar{1 & 2 & 3 \\ 4 & 5 & 6 \\ 1 & 2 & 3} = 0
    \]
  \end{block}

  \begin{block}{Por qué}
    Intercambia los dos renglones iguales. La matriz no cambió, pero
    la regla 2 dice que el signo sí:
    \[
      D = -D
      \quad\Longrightarrow\quad
      D = 0
    \]
  \end{block}
\end{frame}

\begin{frame}{Regla 5: sumar o restar un múltiplo no cambia nada}
  \begin{block}{Ejemplo}
    \[
      \dbar{1 & 2 \\ 3 & 7} = 7 - 6 = 1
      \xrightarrow{\ R_2 - 3R_1\ }
      \dbar{1 & 2 \\ 0 & 1} = 1 - 0 = 1
    \]
  \end{block}

  \[
    \dbar{a & b \\ c \pm \ell a & d \pm \ell b}
    = \dbar{a & b \\ c & d}
      \pm \ell\,\underbrace{\dbar{a & b \\ a & b}}_{0}
    = \dbar{a & b \\ c & d}
  \]


\end{frame}

\begin{frame}{Regla 6: un renglón de ceros da cero}
  \begin{block}{Ejemplo}
    \[
      \dbar{5 & 9 \\ 0 & 0} = 0 - 0 = 0
    \]
  \end{block}

  \begin{block}{Por qué}
    El renglón de ceros es $0$ por cualquier otro renglón. La regla
    3 con $t = 0$ saca ese cero afuera del determinante.
  \end{block}
\end{frame}

\begin{frame}{Regla 7: triangular es el producto de la diagonal}
  \[
    \dbar{d_1 & * & * \\ & \ddots & * \\ 0 & & d_n}
      = d_1 d_2 \cdots d_n
  \]

  \begin{block}{Ejemplo}
    \[
      \dbar{2 & 7 & 1 \\ 0 & 3 & 5 \\ 0 & 0 & 4}
      = (2)(3)(4) = 24
    \]
  \end{block}

  \begin{itemize}
    \item Limpiar con la regla 5, factorizar con la regla 3, cerrar
          con la regla 1.
    \item Si algún $d_i = 0$: renglón de ceros, determinante $0$.
  \end{itemize}
\end{frame}

\begin{frame}{Regla 8: cero si y solo si es singular}
  \[
    \dt{A} = \pm\,(d_1 d_2 \cdots d_n)
    = \pm\,(\text{producto de los pivotes})
  \]

  \begin{block}{Ejemplo: singular}
    El renglón 2 es $2$ veces el renglón 1.
    \[
      \dbar{1 & 2 \\ 2 & 4}
      \xrightarrow{\ R_2 - 2R_1\ }
      \dbar{1 & 2 \\ 0 & 0} = 0
    \]
  \end{block}

  \begin{itemize}
    \item Singular: falta un pivote, aparece un renglón de ceros.
    \item Invertible: los $n$ pivotes son distintos de cero.
  \end{itemize}
\end{frame}

\begin{frame}{Regla 9: el determinante de producto}
  \[
    \dt{(AB)} = (\dt{A})(\dt{B})
  \]

  \begin{block}{Ejemplo}
    \[
      \underbrace{\mat{1 & 2 \\ 3 & 4}}_{\dt{} \,=\, -2}
      \underbrace{\mat{2 & 0 \\ 1 & 3}}_{\dt{} \,=\, 6}
      = \underbrace{\mat{4 & 6 \\ 10 & 12}}_{\dt{} \,=\, -12}
      \qquad
      (-2)(6) = -12
    \]
  \end{block}

  \remate{Con $B = A^{-1}$ sale $(\dt{A})(\dt{A^{-1}}) = 1$.}
\end{frame}

\begin{frame}{Regla 10: determinante de la transpuesta}
  \[
    \dt{(A^{\mathsf{T}})} = \dt{A}
  \]

  \begin{block}{Ejemplo}
    \[
      \dbar{1 & 2 & 3 \\ 0 & 4 & 5 \\ 1 & 0 & 6} = 22
      \qquad
      \dbar{1 & 0 & 1 \\ 2 & 4 & 0 \\ 3 & 5 & 6} = 22
    \]
  \end{block}


\end{frame}

% =====================================================================
%  3. Cálculo por eliminación gaussiana
% =====================================================================
\section{Eliminación gaussiana (pivotes)}

\begin{frame}{El algoritmo práctico}
  \[
    PA = LU
    \quad\Longrightarrow\quad
    (\dt{P})(\dt{A}) = (\dt{L})(\dt{U})
    \quad\Longrightarrow\quad
    \dt{A} = \pm\,(d_1 \cdots d_n)
  \]

  \begin{itemize}
    \item $\dt{L} = 1$: unos en la diagonal.
    \item $\dt{U}$: producto de los pivotes, regla 7.
    \item Signo: paridad de los intercambios.
  \end{itemize}

  \begin{block}{Costo}
    $\sim n^3/3$ operaciones.
  \end{block}
\end{frame}

\begin{frame}{Ejemplo}
  \footnotesize
  \[
    \mat{0 & 1 & 2 \\ 1 & 2 & 3 \\ 2 & 1 & 1}
    \xrightarrow{\ R_1 \leftrightarrow R_2\ }
    \mat{1 & 2 & 3 \\ 0 & 1 & 2 \\ 2 & 1 & 1}
    \xrightarrow{\ R_3 - 2R_1\ }
    \mat{1 & 2 & 3 \\ 0 & 1 & 2 \\ 0 & -3 & -5}
    \xrightarrow{\ R_3 + 3R_2\ }
    \mat{1 & 2 & 3 \\ 0 & 1 & 2 \\ 0 & 0 & 1}
  \]
  \normalsize

  \begin{itemize}
    \item $a_{11} = 0$ no sirve de pivote: hay que intercambiar.
    \item Pivotes: $1$, $1$, $1$. Producto $= 1$.
    \item Un intercambio, número impar, así que $\dt{P} = -1$.
  \end{itemize}

  \[
    \dt{A} = -(1)(1)(1) = -1
  \]
\end{frame}

\begin{frame}{Ejemplo: dónde está $PA = LU$}
  Las tres matrices ya estaban en la cadena de flechas.
  \[
    \underbrace{\mat{0 & 1 & 0 \\ 1 & 0 & 0 \\ 0 & 0 & 1}}_{P}
    \qquad
    \underbrace{\mat{1 & 0 & 0 \\ 0 & 1 & 0 \\ 2 & -3 & 1}}_{L}
    \qquad
    \underbrace{\mat{1 & 2 & 3 \\ 0 & 1 & 2 \\ 0 & 0 & 1}}_{U}
  \]

  \begin{itemize}
    \item $P$: el intercambio $R_1 \leftrightarrow R_2$.
    \item $L$: los multiplicadores, cada uno donde se usó.
    \item $U$: la última matriz de la cadena.
  \end{itemize}

  \begin{alertblock}{$L$ guarda lo que se \emph{resta}}
    La flecha decía $R_3 + 3R_2$, que es $R_3 - (-3)R_2$. En $L$ va
    $-3$, no $+3$.
  \end{alertblock}
\end{frame}

\begin{frame}{Ejemplo: y de ahí el determinante}
  \[
    LU = \mat{1 & 0 & 0 \\ 0 & 1 & 0 \\ 2 & -3 & 1}
         \mat{1 & 2 & 3 \\ 0 & 1 & 2 \\ 0 & 0 & 1}
       = \mat{1 & 2 & 3 \\ 0 & 1 & 2 \\ 2 & 1 & 1}
       = PA
  \]

  \begin{itemize}
    \item $\dt{P} = -1$: un intercambio, regla 2.
    \item $\dt{L} = 1$: unos en la diagonal, regla 7. El $2$ y el
          $-3$ no cuentan.
    \item $\dt{U} = (1)(1)(1) = 1$: triangular, regla 7.
  \end{itemize}

  \[
    (\dt{P})(\dt{A}) = (\dt{L})(\dt{U})
    \quad\Longrightarrow\quad
    (-1)\,\dt{A} = 1
    \quad\Longrightarrow\quad
    \dt{A} = -1
  \]
\end{frame}


% =====================================================================
%  4. La fórmula grande
% =====================================================================
\section{La fórmula grande (permutaciones)}

\begin{frame}{Partir renglones}
  \[
    \dbar{a & b \\ c & d}
      = \dbar{a & 0 \\ c & d} + \dbar{0 & b \\ c & d}
      = \dbar{a & 0 \\ 0 & d} + \dbar{0 & b \\ c & 0}
  \]

  \[
    = ad\,\dbar{1 & 0 \\ 0 & 1} + bc\,\dbar{0 & 1 \\ 1 & 0}
      = ad - bc
  \]

  \remate{Lo que sobrevive son \alert{matrices de permutación}.
          Ellas ponen los signos.}
\end{frame}

\begin{frame}{La fórmula 3 por 3: seis términos}
  \[
    \dbar{a_{11} & a_{12} & a_{13} \\
          a_{21} & a_{22} & a_{23} \\
          a_{31} & a_{32} & a_{33}}
    =
    \begin{aligned}
      &+ a_{11}a_{22}a_{33} + a_{12}a_{23}a_{31}
         + a_{13}a_{21}a_{32} \\
      &- a_{11}a_{23}a_{32} - a_{12}a_{21}a_{33}
         - a_{13}a_{22}a_{31}
    \end{aligned}
  \]

  Cada término se factoriza: el coeficiente por una \alert{matriz de
  permutación}. Ella pone el signo.

  \footnotesize
  \[
    \begin{aligned}
      \dt{A} =\;
        & a_{11}a_{22}a_{33}\dbar{1 & 0 & 0 \\ 0 & 1 & 0 \\ 0 & 0 & 1}
        + a_{12}a_{23}a_{31}\dbar{0 & 1 & 0 \\ 0 & 0 & 1 \\ 1 & 0 & 0}
        + a_{13}a_{21}a_{32}\dbar{0 & 0 & 1 \\ 1 & 0 & 0 \\ 0 & 1 & 0}
        \\[0.4em]
        + & a_{11}a_{23}a_{32}\dbar{1 & 0 & 0 \\ 0 & 0 & 1 \\ 0 & 1 & 0}
        + a_{12}a_{21}a_{33}\dbar{0 & 1 & 0 \\ 1 & 0 & 0 \\ 0 & 0 & 1}
        + a_{13}a_{22}a_{31}\dbar{0 & 0 & 1 \\ 0 & 1 & 0 \\ 1 & 0 & 0}
    \end{aligned}
  \]
  \normalsize



\end{frame}

\begin{frame}{Generalización}
  \begin{block}{Suma sobre las $n!$ permutaciones de columnas}
    \[
      \dt{A} = \sum_{P} (\dt{P})\, a_{1\alpha}a_{2\beta}\cdots
               a_{n\omega}
    \]
  \end{block}

  \begin{itemize}
    \item $(\alpha,\beta,\ldots,\omega)$: los ordenamientos de
          $(1,\ldots,n)$.
    \item $\dt{P} = \pm 1$ según la paridad.
    \item $U$ triangular: solo sobrevive
          $u_{11}u_{22}\cdots u_{nn}$.
  \end{itemize}
\end{frame}

\begin{frame}{Pq no se usa}
  \begin{center}
    \begin{tabular}{r r r}
      \toprule
      $n$ & términos $n!$ & operaciones $\approx n^3/3$ \\
      \midrule
       3 & 6                 & 9 \\
       5 & 120               & 42 \\
      10 & 3\,628\,800       & 333 \\
      11 & 39\,916\,800      & 444 \\
      20 & $2.4\times10^{18}$ & 2\,667 \\
      \bottomrule
    \end{tabular}
  \end{center}

  \begin{alertblock}{El truco de las diagonales}
    Sirve solo para $n = 3$. Con $n = 4$ da $8$ productos y hacen
    falta $24$.
  \end{alertblock}
\end{frame}

% =====================================================================
%  5. Cofactores
% =====================================================================
\section{Expansión por cofactores}

\begin{frame}{Algoritmo}
  Tacha el renglón $1$ y la columna $j$: sobra el menor $M_{1j}$.

  \begin{center}
    \scriptsize
    \begin{tikzpicture}[every node/.style={font=\scriptsize}]
      \node (a1) at (0, 1.25)
        {$\mat{\ta{a_{11}} & \ta{a_{12}} & \ta{a_{13}} \\
               \ta{a_{21}} & a_{22} & a_{23} \\
               \ta{a_{31}} & a_{32} & a_{33}}$};
      \node (a2) at (0, 0)
        {$\mat{\ta{a_{11}} & \ta{a_{12}} & \ta{a_{13}} \\
               a_{21} & \ta{a_{22}} & a_{23} \\
               a_{31} & \ta{a_{32}} & a_{33}}$};
      \node (a3) at (0,-1.25)
        {$\mat{\ta{a_{11}} & \ta{a_{12}} & \ta{a_{13}} \\
               a_{21} & a_{22} & \ta{a_{23}} \\
               a_{31} & a_{32} & \ta{a_{33}}}$};

      \node (m1) at (3.1, 1.25)
        {$M_{11} = \dbar{a_{22} & a_{23} \\ a_{32} & a_{33}}$};
      \node (m2) at (3.1, 0)
        {$M_{12} = \dbar{a_{21} & a_{23} \\ a_{31} & a_{33}}$};
      \node (m3) at (3.1,-1.25)
        {$M_{13} = \dbar{a_{21} & a_{22} \\ a_{31} & a_{32}}$};

      \draw[->] (a1) -- (m1);
      \draw[->] (a2) -- (m2);
      \draw[->] (a3) -- (m3);

      \node (r1) at (6.8, 2.0) {$a_{33}$};
      \node (r2) at (6.8, 1.55) {$a_{32}$};
      \node (r3) at (6.8, 0.95) {$a_{23}$};
      \node (r4) at (6.8, 0.5) {$a_{22}$};
      \draw[->,acento] (m1) -- (r1);
      \draw[->,acento] (m1) -- (r2);
      \draw[->,gray]   (m1) -- (r3);
      \draw[->,gray]   (m1) -- (r4);
    \end{tikzpicture}
  \end{center}

  \remate{\[
    C_{1j} = (-1)^{1+j}\,M_{1j}
    \qquad\Longrightarrow\qquad
    C_{11} = +M_{11}
    \quad
    C_{12} = -M_{12}
    \quad
    C_{13} = +M_{13}
  \]}

\end{frame}



\begin{frame}{Cofactores y el tablero de ajedrez}
  \begin{block}{Tacha el renglón $i$ y la columna $j$}
    \[
      C_{ij} = (-1)^{i+j}\,\dt{M_{ij}}
      \qquad
      M_{ij} \text{ de tamaño } n-1
    \]
  \end{block}

  \[
    (-1)^{i+j} =
    \mat{+ & - & + & - \\
         - & + & - & + \\
         + & - & + & - \\
         - & + & - & +}
  \]
\end{frame}

\begin{frame}{La expansión}
  \[
    \dt{A} = a_{i1}C_{i1} + \cdots + a_{in}C_{in}
    \qquad \text{(renglón $i$)}
  \]
  \[
    \dt{A} = a_{1j}C_{1j} + \cdots + a_{nj}C_{nj}
    \qquad \text{(columna $j$)}
  \]

  \begin{itemize}
    \item Orden $n$ a orden $n-1$. Recursiva hasta $\dt{[a]} = a$.
  \end{itemize}
\end{frame}


% =====================================================================
%  6. Cierre
% =====================================================================
\section{Cómo elegir el método}

\begin{frame}{Qué método y cuándo}
  \begin{center}
    \small
    \begin{tabular}{@{}l l l@{}}
      \toprule
      \textbf{Método} & \textbf{Costo} & \textbf{Cuándo usarlo} \\
      \midrule
      Eliminación gaussiana & $\sim n^3/3$ & Matriz numérica y
                                            general \\
      Fórmula grande & $n!$ & Entradas explícitas, o $n \leq 3$ \\
      Cofactores & $n!$ peor caso & Dispersa, de banda o
                                    simbólica \\
      \bottomrule
    \end{tabular}
  \end{center}

  \remate{Eliminación para calcular. Cofactores cuando hay muchos
          ceros. La fórmula grande para demostrar.}
\end{frame}

\begin{frame}{Ideas clave}
  \begin{enumerate}
    \item Definición: $\dt{I} = 1$, cambio de signo, linealidad por
          renglón.
    \item $\dt{A} = \pm(\text{producto de los pivotes})$.
    \item $\dt{A} = 0 \iff A$ no es invertible.
    \item $\dt{(AB)} = (\dt{A})(\dt{B})$ y
          $\dt{(A^{\mathsf{T}})} = \dt{A}$.
    \item Fórmula grande: $n!$ productos con signo.
    \item Cofactores: orden $n$ a orden $n-1$.
  \end{enumerate}
\end{frame}

\begin{frame}{Referencia}
  \begin{block}{Fuentes}
    Gilbert Strang, \emph{Introduction to Linear Algebra}, \\
    capítulo 5: Determinants. 
  \end{block}

\end{frame}

% =====================================================================
%  7. Ejercicios
% =====================================================================
\section{Ejercicios}

\begin{frame}{Ejercicio 1: una matriz 2 por 2}
  \[
    A = \mat{3 & 5 \\ 2 & 4}
  \]

  \begin{block}{Por la formula directa}
    \[
      \dt{A} = (3)(4) - (5)(2) = 12 - 10 = 2
    \]
  \end{block}

  \begin{block}{Por pivotes, para comprobar}
    \[
      \mat{3 & 5 \\ 2 & 4}
      \xrightarrow{\ R_2 - \frac{2}{3}R_1\ }
      \mat{3 & 5 \\ 0 & \tfrac{2}{3}}
      \qquad
      (3)\left(\tfrac{2}{3}\right) = 2
    \]
  \end{block}
\end{frame}

\begin{frame}{Ejercicio 2: una matriz 3 por 3}
  \[
    A = \mat{1 & 2 & 3 \\ 0 & 4 & 5 \\ 1 & 0 & 6}
  \]

  \begin{block}{Cofactores bajando por la primera columna}
    Entradas $1$, $0$, $1$. El cero mata un termino completo.
    \[
      \dt{A}
        = 1\,\dbar{4 & 5 \\ 0 & 6}
        + 1\,\dbar{2 & 3 \\ 4 & 5}
        = (24) + (10 - 12)
        = 22
    \]
  \end{block}

  \begin{alertblock}{Los signos}
    $C_{11}$ y $C_{31}$ son los dos positivos: $(-1)^{1+1}$ y
    $(-1)^{3+1}$.
  \end{alertblock}
\end{frame}

\begin{frame}{Ejercicio 2: la regla de Sarrus}
  Copia las dos primeras columnas a la derecha y multiplica sobre las
  diagonales.

  \begin{center}
    \begin{tikzpicture}[baseline]
      \matrix (m) [matrix of math nodes, column sep=1.1em,
                   row sep=0.7em,
                   nodes={minimum width=1.1em, inner sep=2pt,
                          fill=white, rounded corners=1pt},
                   ampersand replacement=\&] {
        1 \& 2 \& 3 \& 1 \& 2 \\
        0 \& 4 \& 5 \& 0 \& 4 \\
        1 \& 0 \& 6 \& 1 \& 0 \\
      };
      % Las diagonales van al fondo para no tapar los digitos.
      \begin{scope}[on background layer]
        \draw[gray!50] ($(m-1-3)!0.5!(m-1-4) + (0,1em)$)
                    -- ($(m-3-3)!0.5!(m-3-4) - (0,1em)$);
        \draw[acento,line width=1.1pt]          (m-1-1) -- (m-3-3);
        \draw[acento,line width=1.1pt]          (m-1-2) -- (m-3-4);
        \draw[acento,line width=1.1pt]          (m-1-3) -- (m-3-5);
        \draw[orange!70!black,line width=1.1pt] (m-3-1) -- (m-1-3);
        \draw[orange!70!black,line width=1.1pt] (m-3-2) -- (m-1-4);
        \draw[orange!70!black,line width=1.1pt] (m-3-3) -- (m-1-5);
      \end{scope}
    \end{tikzpicture}
  \end{center}

  \[
    \dt{A} =
    \underbrace{(24) + (10) + (0)}_{\text{\color{acento} bajando}}
    - \underbrace{(12) + (0) + (0)}
      _{\text{\color{orange!70!black} subiendo}}
    = 34 - 12 = 22
  \]

  \remate{Mismo 22 que dieron los cofactores.
          Solo sirve para $3$ por $3$.}
\end{frame}

\begin{frame}{Ejercicio 3: una matriz 4 por 4}
  \[
    A = \mat{1 & 2 & 0 & 1 \\
             2 & 3 & 1 & 0 \\
             0 & 1 & 2 & 1 \\
             1 & 0 & 1 & 2}
  \]

  Con $4! = 24$ terminos la formula grande no es opcion y la primera
  columna no tiene suficientes ceros para cofactores.

  \remate{Se elimina.}
\end{frame}

\begin{frame}{Ejercicio 3: el procedimiento}
  \footnotesize
  \[
    \mat{1 & 2 & 0 & 1 \\
         2 & 3 & 1 & 0 \\
         0 & 1 & 2 & 1 \\
         1 & 0 & 1 & 2}
    \xrightarrow[R_4 - R_1]{R_2 - 2R_1}
    \mat{1 & 2 & 0 & 1 \\
         0 & -1 & 1 & -2 \\
         0 & 1 & 2 & 1 \\
         0 & -2 & 1 & 1}
    \xrightarrow[R_4 - 2R_2]{R_3 + R_2}
    \mat{1 & 2 & 0 & 1 \\
         0 & -1 & 1 & -2 \\
         0 & 0 & 3 & -1 \\
         0 & 0 & -1 & 5}
  \]

  \[
    \xrightarrow{\ R_4 + \tfrac{1}{3}R_3\ }
    \mat{1 & 2 & 0 & 1 \\
         0 & -1 & 1 & -2 \\
         0 & 0 & 3 & -1 \\
         0 & 0 & 0 & \tfrac{14}{3}}
    \qquad
    \dt{A} = (1)(-1)(3)\left(\tfrac{14}{3}\right) = -14
  \]

  \normalsize
  \remate{Sin intercambios, el signo es $+$. El determinante sale
          negativo por el pivote $-1$, no por la permutación.}
\end{frame}

\begin{frame}{Ejercicio 4: para resolver}
  \[
    B = \mat{0 & 1 & 2 & 0 \\
             1 & 0 & 3 & 1 \\
             2 & 1 & 0 & 2 \\
             0 & 2 & 1 & 1}
  \]

  \remate{Calcular $\dt{B}$.}

\end{frame}

\begin{frame}[standout]
  Preguntas
\end{frame}

% =====================================================================
\end{document}
% =====================================================================
%########################## END OF DETERMINANTES.TEX ###########################
%###############################################################################
